\chapter*{Conclusion générale}
\addcontentsline{toc}{chapter}{Conclusion générale}
\label{ch:conclusion}

Au terme de ce rapport, consacré à la conception et à l'évaluation
d'une méthode de détection d'anomalies dans le trafic des réseaux
cellulaires 4G/5G, il convient de revenir sur les apports de ce
travail, ses limites, et les perspectives qu'il ouvre. La
problématique posée en introduction était double~: comment distinguer
une variation normale du trafic mobile -- heures de pointe,
week-ends, saisonnalité -- d'un comportement réellement atypique, à
l'échelle d'une métropole entière et de plusieurs canaux de
communication, sans disposer d'anomalies déjà étiquetées pour
entraîner un modèle supervisé~? Pour y répondre, ce travail s'est
appuyé sur le framework académique
DiNATrAX~\parencite{maudoux2024dinatrax}, en l'appliquant à des
données réelles de trafic Internet, SMS et appels de la ville de
Milan.

\section*{Synthèse des contributions}

Les résultats obtenus répondent, un à un, aux objectifs spécifiques
fixés au chapitre~\ref{chap:contexte}. Le profil horaire de trafic de
chaque grille a été encodé sous une forme compacte et comparable --
une signature de 24 lettres, dont les six niveaux sont calibrés sur
les percentiles réels du trafic plutôt que sur des seuils ronds
arbitraires, ce qui rééquilibre le pouvoir discriminant de la
signature (section~\ref{sec:comparaison-encodage}). Cette signature,
comparée à un profil de référence par distance de
Damerau-Levenshtein puis convertie en score z, permet d'isoler, sans
aucune donnée étiquetée, un ensemble restreint de journées atypiques
par grille~: 2\,484 sur le canal Internet, 4\,755 sur SMS et 6\,468
sur les appels (tableau~\ref{tab:resultats-canal}), avec des taux de
correspondance calendaire allant de 86,8~\% à 97,6~\% selon le canal.
Appliquée indépendamment aux trois canaux, la détection confirme
surtout l'intérêt d'un croisement multi-canal~: une anomalie signalée
simultanément par au moins deux canaux atteint 99,2~\% de
correspondance calendaire, contre 86,8 à 97,6~\% pour un canal isolé
(section~\ref{sec:multicanal}) -- la preuve la plus nette que la
méthode retenue produit un signal fiable plutôt qu'un simple artefact
statistique.

Un raffinement par fenêtre glissante a en outre permis de réduire la
latence de détection, initialement bornée à 24 heures par la
granularité journalière du bloc de référence, jusqu'à un gain de
11 heures sur le cas étudié (grille 1421, 1\textsuperscript{er} janvier
2014, tableau~\ref{tab:latence}), sans modifier le cœur de la méthode
-- uniquement la fréquence de réévaluation. Ce raffinement montre
qu'une contrainte opérationnelle de supervision en quasi temps réel
peut être satisfaite à moindre coût de calcul. Enfin, l'ensemble de
ces résultats est explorable via un tableau de bord 3D interactif
(deck.gl, MapLibre GL JS), qui rend directement exploitables les
anomalies détectées pour un usage opérationnel de supervision réseau.

\section*{Limites}

Ce travail présente néanmoins plusieurs limites, déjà en partie
identifiées au chapitre~\ref{chap:contexte} et confirmées par les
résultats du chapitre~\ref{chap:conception}. D'abord, les données
utilisées portent sur une seule ville et une période de deux mois
(novembre-décembre 2013)~: les seuils de nomenclature, calibrés sur
les percentiles réels de ce jeu de données précis, ne sont donc pas
directement transposables à un autre contexte géographique, temporel,
ou à un opérateur différent, sans nouvelle calibration. Ensuite, la
validation des anomalies détectées repose sur un calendrier de repère
construit manuellement (jours fériés, matchs à San Siro), qui reste
partiel par construction~: les anomalies sans explication calendaire
identifiée (13,2~\% à 2,4~\% selon le canal, moins de 1~\% pour les
anomalies confirmées) ne sont pas nécessairement des faux positifs,
mais en l'absence de vérité terrain fournie par un opérateur réseau,
il n'est pas possible de trancher, au cas par cas, entre un événement
réel non répertorié et un signal résiduel sans cause identifiable.
Enfin, le jeu de données ne distingue pas les usages 4G des usages 5G
à proprement parler, la technologie 5G n'étant pas encore déployée
sur la période couverte~; l'approche reste conceptuellement
transposable à des données de trafic 5G actuelles, mais cette
transposition n'a pas été testée dans ce travail.

\section*{Perspectives}

Ces limites ouvrent plusieurs pistes de prolongement concrètes.
Premièrement, un rapprochement avec les équipes d'exploitation d'un
opérateur réseau permettrait de confronter les anomalies détectées à
un historique réel d'incidents (pannes, maintenance, saturation), et
de remplacer le calendrier de repère -- un proxy utile mais
imparfait -- par une véritable vérité terrain, seule capable de
mesurer précision et rappel effectifs de la méthode. Deuxièmement, la
même méthode pourrait être appliquée à d'autres villes, sur une
période plus longue et sur des données 5G actuelles, afin de vérifier
la stabilité de l'approche et de déterminer si les seuils doivent
être recalibrés à chaque déploiement ou si des valeurs par défaut
raisonnables peuvent être dégagées. Troisièmement, la détection par
fenêtre glissante, déjà conçue pour ne coûter qu'une comparaison de
deux séquences de 24 caractères par évaluation horaire
(section~\ref{sec:fenetre-glissante}), ouvre la voie à un système
d'alerte fonctionnant directement sur un flux de trafic en continu,
plutôt que sur une analyse a posteriori de fichiers journaliers.
Quatrièmement, le calendrier de repère pourrait être enrichi d'autres
sources d'événements locaux (travaux, incidents météorologiques,
actualité de la ville), afin de réduire la part d'anomalies
résiduelles laissées sans explication. Enfin, une fois une vérité
terrain disponible grâce aux pistes précédentes, il deviendrait
possible d'entraîner des modèles supervisés ou semi-supervisés et de
les comparer à la méthode non supervisée développée ici, qui
conserverait alors son rôle de référence robuste, applicable sans
délai à un nouveau déploiement dépourvu de données étiquetées.

En définitive, ce travail montre qu'une méthode de détection
d'anomalies simple dans son principe -- encoder, comparer, seuiller --
peut, appliquée avec rigueur à des données réelles et validée par un
croisement multi-canal, produire un signal suffisamment fiable pour
répondre à un enjeu opérationnel concret de supervision des réseaux
mobiles 4G/5G, sans jamais nécessiter de données étiquetées au
préalable.
